\documentclass[11pt]{article}
\usepackage[utf8]{inputenc}
\usepackage[T1]{fontenc}
\usepackage{lmodern}
\usepackage[margin=1in]{geometry}
\usepackage{amsmath,amssymb,mathtools}
\usepackage{graphicx}
\usepackage{grffile}
\usepackage{float}
\usepackage{hyperref}
\usepackage{caption}
\usepackage{parskip}
\usepackage{enumitem}
\usepackage{xurl}
\setlength{\parskip}{0.65em}
\setlength{\parindent}{0pt}
\begin{document}
\title{Daily Research Digest --- 2026-04-23}
\author{}
\date{}
\maketitle
\textbf{Topics:} galaxy clusters, galaxies, gravitational lensing, dark matter.
\section*{Synthesis}
Here's a summary of the key papers and their topics based on the provided abstracts:
\subsection*{Dark Matter:}
1. **"Purely Quadratic Non-Gaussianity from Tachyonic Instability: Primordial Black Holes and Scalar-Induced Gravitational Waves"**
\begin{itemize}
\item Investigates primordial black hole (PBH) formation in scenarios with purely quadratic non-Gaussian curvature perturbations arising from tachyonic instability. Highlights the role of spectral width and correlation coefficients in PBH abundance suppression, while also discussing gravitational-wave signals.
\end{itemize}
2. **"Photon and neutrino fluxes from spheroidal dwarf galaxies in a decaying DM model"**
\begin{itemize}
\item Examines indirect detection signatures (gamma rays and neutrinos) from decaying dark matter models in the context of Milky Way and 14 dwarf spheroidal galaxies, with predictions for specific experiments.
\end{itemize}
3. **"Unraveling Chemical Enrichment in Extreme Emission-Line Galaxies: A Multi-Element Bayesian View of Bursty Star Formation and Galaxy Evolution in DESI"**
\begin{itemize}
\item Analyzes chemical enrichment and star formation histories in extreme emission-line galaxies (EELGs) using a Bayesian model, finding short depletion timescales and large mass-loading factors indicative of burst-driven gas cycling.
\end{itemize}
\subsection*{Galaxies:}
1. **"MAUVE-MUSE: Ionization and Kinematic Signatures of Environmental Effects on Virgo Cluster Disks"**
\begin{itemize}
\item Studies the interstellar medium in 40 galaxies within the Virgo cluster, finding elevated ionization line ratios compared to field disks and significant kinematic broadening consistent with diffuse ionized gas (DIG) and intracluster interactions.
\end{itemize}
2. **"Cold molecular gas distribution and kinematics in the low-metallicity dusty starburst of Mrk 996 resolved with ALMA"**
\begin{itemize}
\item Explores cold molecular gas properties in a nearby, metal-poor Wolf-Rayet galaxy using ALMA observations. Detects compact CO clouds offset from the nuclear super star cluster.
\end{itemize}
3. **"Unraveling Chemical Enrichment in Extreme Emission-Line Galaxies: A Multi-Element Bayesian View of Bursty Star Formation and Galaxy Evolution in DESI"**
\begin{itemize}
\item Investigates chemical enrichment processes in extreme emission-line galaxies (EELGs) using a Bayesian model, highlighting the influence of star formation efficiency, outflows, and inflows.
\end{itemize}
4. **"Photon and neutrino fluxes from spheroidal dwarf galaxies in a decaying DM model"**
\begin{itemize}
\item Also relevant to galaxy evolution through indirect detection studies on dark matter decay products in dwarf galaxies.
\end{itemize}
\subsection*{Astrophysics \& Cosmology:}
1. **"Unraveling Chemical Enrichment in Extreme Emission-Line Galaxies: A Multi-Element Bayesian View of Bursty Star Formation and Galaxy Evolution in DESI"**
\begin{itemize}
\item Also touches upon star formation and galaxy evolution through chemical enrichment studies using a Bayesian framework.
\end{itemize}
2. **"Cold molecular gas distribution and kinematics in the low-metallicity dusty starburst of Mrk 996 resolved with ALMA"**
\begin{itemize}
\item Provides insights into cold molecular gas properties and starbursts in metal-poor galaxies, contributing to our understanding of galaxy evolution and star formation processes.
\end{itemize}
3. **"Unraveling Chemical Enrichment in Extreme Emission-Line Galaxies: A Multi-Element Bayesian View of Bursty Star Formation and Galaxy Evolution in DESI"**
\begin{itemize}
\item Further insights into the chemical enrichment and star formation processes that shape galaxies.
\end{itemize}
\subsection*{Instrumentation \& Data Analysis:}
1. **"Neutron Capture Elements in Metal-Poor Stars from SDSS-V: Implications for Nucleosynthesis and Galactic Chemical Evolution"**
\begin{itemize}
\item Studies neutron capture elements in metal-poor stars using data from SDSS-V, contributing to understanding nucleosynthesis and galactic chemical evolution.
\end{itemize}
2. **"Unraveling Chemical Enrichment in Extreme Emission-Line Galaxies: A Multi-Element Bayesian View of Bursty Star Formation and Galaxy Evolution in DESI"**
\begin{itemize}
\item Uses a Bayesian model for multi-element analysis, providing insights into star formation histories through detailed data analysis techniques.
\end{itemize}
\subsection*{Summary:}
The papers cover a range of topics including dark matter decay and its detection via photons and neutrinos from dwarf galaxies, the role of quadratic non-Gaussianity in PBH formation, chemical enrichment processes in extreme emission-line galaxies (EELGs), cold molecular gas properties in metal-poor starbursts, and environmental effects on galactic disks within clusters like Virgo. The methodologies span multi-element Bayesian modeling, detailed spectral line analysis using ALMA data, and indirect detection techniques for dark matter studies.
\newpage
\section*{Paper Summaries and Figures}
\subsection*{1. Self-regulated galaxy evolution within a self-consistently varying galaxy-wide IMF}
\textbf{Focus:} galaxies\\
\textbf{Authors:} Lukas Hof, Pavel Kroupa, Gerhard Hensler, Jan Pflamm-Altenburg\\
\textbf{Published:} 2026-04-22T17:59:59+00:00\\
\textbf{PDF:} \href{https://arxiv.org/pdf/2604.20843v1}{https://arxiv.org/pdf/2604.20843v1}\\
\textbf{Summary:}
Semi-analytical evolution models of galaxies are a useful and computationally inexpensive tool for fast assessment of individual properties and their evolution. In this work, specifically the influence of a metallicity and star-formation rate (SFR) dependent galaxy-wide stellar initial mass function (IGIMF) on the self-regulation of star-formation in a galaxy is of interest. All models -- both non-varying gwIMFs and the IGIMF -- reproduce reasonable gas fractions, gas depletion timescales and the main sequence of star-forming galaxies. However, only the IGIMF model accurately predicts the mass-metallicity relation and provides a more comprehensive description of quenched elliptical galaxies. For massive ellipticals all models suggest the need for an additional gas heating source to reach a quenched state. Using a different stellar yield table in the IGIMF model does not significantly affect the results. In all models, the galaxies evolve self-regulated, determined by the accretion rate. The self-regulated constancy of the SFR reflects the constant SFRs of nearby star-forming galaxies. The specific gas-accretion rate of all galaxies appears to be comparable to the Hubble constant. The inclusion of outflows improves the results for the canonical gwIMF model, but not significantly, while for the IGIMF model it has no significant impact.
\newpage
\subsection*{2. Search for Axion Like Particles produced via the Primakoff process at COMPASS}
\textbf{Focus:} dark matter\\
\textbf{Authors:} Mehran Dehpour\\
\textbf{Published:} 2026-04-22T16:20:45+00:00\\
\textbf{PDF:} \href{https://arxiv.org/pdf/2604.20734v1}{https://arxiv.org/pdf/2604.20734v1}\\
\textbf{Summary:}
Axion-Like Particles (ALPs) are well-motivated candidates for dark matter and potential mediators to the dark sector. We present a search for ALPs coupled to photons, based on a reinterpretation of COMPASS data. Using the 2009 dataset consisting of $190~\text{GeV}$ $\pi ^-$ and $\mu ^-$ beams impinging on a fixed nickel target, we investigate the Primakoff production of ALPs. Due to the high beam energy, ALPs in the MeV mass range are produced with a significant Lorentz boost, leading to strongly collimated decay photons. Consequently, these photons are not spatially resolved by the electromagnetic calorimeter and are instead reconstructed as a single merged cluster. This signature mimics the single-photon signal of Standard Model Primakoff Compton scattering, which was the primary focus of the original COMPASS analysis. By quantifying this potential ALP contamination in the Compton scattering sample, we derive exclusion limits on the ALP-photon coupling $g_{a\gamma \gamma }$ in the mass range $0.2 \lesssim m_a \lesssim 600~\text{MeV}$. Our results exclude couplings $g_{a\gamma \gamma } \gtrsim 10^{-1}~\text{GeV}^{-1}$ at 95\% C.L., providing independent constraints on the parameter space that bridges beam dump experiments and high energy colliders. While current collider-based limits remain more stringent, this work establishes a novel reinterpretation framework and provides a baseline for future studies of resolved diphoton states in complementary kinematic regimes, such as Primakoff $\pi ^0$ production.
\newpage
\subsection*{3. Dilaton-Induced Resonant Production of Ultralight Vector Dark Matter}
\textbf{Focus:} dark matter\\
\textbf{Authors:} Imtiaz Khan, G. Mustafa, Jehanzad Zafar, Farruh Atamurotov, Ahmadjon Abdujabbarov, Chengxun Yuan\\
\textbf{Published:} 2026-04-22T12:14:09+00:00\\
\textbf{PDF:} \href{https://arxiv.org/pdf/2604.20481v1}{https://arxiv.org/pdf/2604.20481v1}\\
\textbf{Summary:}
We study the resonant production of ultralight vector dark matter from an oscillating spectator scalar $\phi $ coupled to a massive vector $A_\mu $ through a dilatonic kinetic function $\W(\phi )$. The mechanism contains a narrow branch near $\mA/\mphi\simeq 1/2$. Assuming that the spectator remains subdominant until oscillations begin, we derive the onset fraction $r_i=\Phii^2/(6\Mpl^2)$ and show that, in the linear regime and for a background with constant equation-of-state parameter $w_b$, the growth-to-Hubble ratio scales as $\mu /H\propto a^{3w_b/2}$. Combined with the tuned-branch abundance estimate, this implies $m_{\gamma '}\propto r_i^{-2}$ for the relic dark-photon mass. In particular, the interval $m_{\gamma '}\sim10^{-20}$--$10^{-18}\,\mathrm{eV}$ maps to $r_i\sim10^{-4}$--$10^{-5}$, corresponding to a radiation-dominated onset with a subdominant spectator, while for $M=10^{17}\,\mathrm{GeV}$ the perturbative range $ε_i=0.1$--$1$ gives $m_{\gamma '}\sim10^{-17}$--$10^{-21}\,\mathrm{eV}$. We also derive the polarization-resolved quadratic action in an FLRW background and formulate ultraviolet consistency conditions for both Stückelberg and Higgs completions, including perturbative control of the kinetic modulation, dark-Higgs decoupling, and symmetry-restoration bounds from vector backreaction.
\subsubsection*{Selected figures}
\begin{figure}[H]
\centering
\includegraphics[width=\linewidth,height=0.42\textheight,keepaspectratio]{/home/kf/research-assistant/data/cache/figures/http-arxiv.org-abs-2604.20481v1/figure\_1.png}
\captionsetup{labelformat=empty}\caption{FIG. 3: Floquet growth rates for the tuned branch, mA/mϕ = 1/2, in the transverse and longitudinal sectors
as functions of the dimensionless momentum k/mϕ and amplitude ϕ0/M. The transverse instability occupies a
broader momentum range and develops a secondary higher-k tongue, whereas the longitudin}
\end{figure}
\newpage
\subsection*{4. How Invisible: Regressing The Key Model Parameter for Semi-visible Jet Searches}
\textbf{Focus:} dark matter\\
\textbf{Authors:} Yin Li, Bingxuan Liu, Jianbin Wang, Jiaqi Xie, Kairong Xu, Ruihan Ye, Zihuan Huang\\
\textbf{Published:} 2026-04-22T11:25:37+00:00\\
\textbf{PDF:} \href{https://arxiv.org/pdf/2604.20456v1}{https://arxiv.org/pdf/2604.20456v1}\\
\textbf{Summary:}
Semi-visible jets (SVJs) provide a characteristic collider signature of strongly interacting dark sectors, in which the key model parameter $r_{\mathrm{inv}}$ controls the fraction of dark hadrons decaying to dark matter candidates. In this work, a regression model is developed to reconstruct $r_{\mathrm{inv}}$ in SVJ events produced in association with an energetic photon. The model uses information from high-level physics objects only, and the training procedure is optimized to ensure applicability. The performance is found to be robust against varying signal parameters and $r_{\mathrm{inv}}$ can be reconstructed at a much higher precision, compared to previously developed analytical method. It offers a new approach to conduct SVJ searches that can potentially unify both $s$-channel and $t$-channel productions, enhancing the sensitivities.
\newpage
\subsection*{5. The Evolution of the SFR-M\_* relation at 0.1<z<4: Environmental and Morphological Dependencies}
\textbf{Focus:} galaxies\\
\textbf{Authors:} Kaimin He, Ke Shi, Jun Toshikawa, Xianzhong Zheng, Xiaopeng You\\
\textbf{Published:} 2026-04-22T10:29:47+00:00\\
\textbf{PDF:} \href{https://arxiv.org/pdf/2604.20411v1}{https://arxiv.org/pdf/2604.20411v1}\\
\textbf{Summary:}
We present a comprehensive study of the relationship between star formation rate (SFR) and stellar mass (M\_*) from z = 0.1 to z = 4 using a mass-complete sample of approximately 290,000 galaxies from the COSMOS2020 catalog. We find that the SFR-M\_* relation exhibits a pronounced high-mass decline that becomes increasingly evident at lower redshifts. Examining environmental and morphological dependencies, we find strikingly different patterns. For all galaxies, we find galaxies in high-density environments exhibit suppressed star formation rates at z < 1 especially at high-mass end, while for star-forming galaxies no apparent environmental effect is found at all redshifts. In contrast, galaxy morphology exerts strong influence on the SFR-M\_* relation at z < 2, in a sense that early-type galaxies exhibit systematically lower star formation rates at fixed mass compared to spirals and irregulars, with this trend persisting even within the star-forming population. These results suggest that internal structural properties (bulge components in particular) continuously regulate star formation efficiency independently of whether galaxies are classified as active or quiescent, whereas external environmental processes primarily serve as rapid quenching mechanisms that increase the fraction of quiescent galaxies at low redshifts. We attribute the observed high-mass decline of the SFR-M\_* relation to COSMOS2020's superior capability for detecting massive star-forming galaxies undergoing "morphological quenching" processes.
\newpage
\subsection*{6. Prospects of boosted magnetic dipole inelastic fermion dark matter at ILC-BDX}
\textbf{Focus:} dark matter\\
\textbf{Authors:} I. V. Voronchikhin, D. V. Kirpichnikov\\
\textbf{Published:} 2026-04-22T09:39:06+00:00\\
\textbf{PDF:} \href{https://arxiv.org/pdf/2604.20385v1}{https://arxiv.org/pdf/2604.20385v1}\\
\textbf{Summary:}
In this work, we investigate the projected sensitivity of the Beam-Dump eXperiment at the International Linear Collider (ILC-BDX) to inelastic fermionic dark matter coupled to the Standard Model photon through an off-diagonal magnetic dipole operator. We compute the production rate of dark matter states in the bremsstrahlung like process $e^- N \to e^- N \gamma ^* (\to χ_{1} \barχ_0)$, induced by the scattering of high-energy electrons on target nuclei. The resulting boosted dark matter fluxes are then propagated to the detector, where the signal events arise from scattering off detector electrons. The projected exclusion limits are derived using the expected numbers of electrons on target (this implies a typical rate of $4.0~\times~10^{21}/\mbox{year}$) corresponding to 1 year and 10 years of data taking. To characterize the impact of inelasticity, we consider two benchmark relative mass splittings, $\Delta =0.05$ and $\Delta =0.001$, motivated by thermal dark matter scenarios. Our results show that ILC-BDX can probe inelastic magnetic-dipole dark matter over a phenomenologically relevant region of parameter space.
\subsubsection*{Selected figures}
\begin{figure}[H]
\centering
\includegraphics[width=\linewidth,height=0.42\textheight,keepaspectratio]{/home/kf/research-assistant/data/cache/figures/http-arxiv.org-abs-2604.20385v1/figure\_1.png}
\captionsetup{labelformat=empty}\caption{FIG. 3. Left panel: The expected limits at 95 \% CL on magn
experiment as a function of the lightest DM state, χ0 for mas
the ILC-BDX experiment are shown as red and purple lines fo
while the solid and dashed lines denote 1 year and 10 years
LEP [35, 36] and CHARM II [37, 38] experiments as the gray
}
\end{figure}
\newpage
\subsection*{7. The core of the problem: Physical limits of the core-Sérsic model}
\textbf{Focus:} galaxies\\
\textbf{Authors:} Maarten Baes\\
\textbf{Published:} 2026-04-22T08:35:55+00:00\\
\textbf{PDF:} \href{https://arxiv.org/pdf/2604.20344v1}{https://arxiv.org/pdf/2604.20344v1}\\
\textbf{Summary:}
The core-Sérsic model is the standard tool for describing partially depleted stellar cores in massive early-type galaxies, yet its physical admissibility has rarely been examined. Using numerical deprojections, we show that many formally allowed parameter combinations cannot represent realistic stellar systems: sharp transitions between the inner power-law core and the outer Sérsic profile (large $\alpha $) always generate non-monotonic intrinsic density profiles. We identify, for each set of structural parameters $(\gamma , m, R_{\text{e}}/R_{\text{b}})$, a critical transition parameter, $\alpha _{\text{crit}}$, above which monotonicity is violated. This threshold systematically depends on the core slope and Sérsic index, implying that a fraction of the commonly used parameter space, including the widely adopted sharp-transition limit $\alpha \rightarrow\infty$, is physically ruled out. These constraints have important consequences for measuring core sizes and mass deficits in massive ellipticals, for constructing dynamical models, and for comparing observations with simulations of supermassive black hole binary evolution.
\newpage
\subsection*{8. Photometric Super-Resolution for Improving Galaxy Morphological Measurements using Conditional Generative Adversarial Networks}
\textbf{Focus:} galaxies, gravitational lensing\\
\textbf{Authors:} Samuel Kahn, Ryan Hausen, Hubert Bretonnière, Nicole Drakos, Brant Robertson\\
\textbf{Published:} 2026-04-22T05:23:02+00:00\\
\textbf{PDF:} \href{https://arxiv.org/pdf/2604.20195v1}{https://arxiv.org/pdf/2604.20195v1}\\
\textbf{Summary:}
The measurement of galaxy morphological parameters from astronomical images features in a wide range of modern analyses, including galaxy evolution and cosmological weak lensing studies. The precision and accuracy of morphological parameter estimation can be influenced by several key factors. The effective seeing of the image, summarized by the point spread function (PSF), limits how galaxy features or light profiles are resolved. The pixel scale of the detector also influences the resolution and the amount of statistical information available for a given object. The depth of the observations determines the signal-to-noise ratio of the image. Improving each of these factors is very costly, either in terms of detector upgrades, observatory design, or observing time. Here, we develop a conditional generative adversarial network, called Neo, trained to transform existing ground-based images into sharper, finer-scale images comparable to space-based image quality. We demonstrate that Neo improves the accuracy of measured morphological parameters by factors of $2$-$10$ when trained to translate Subaru Hyper Suprime-Camera (HSC) images to approximate Hubble Space Telescope (HST) data. Neo is designed for applicability to ongoing, large-scale surveys such as the Legacy Survey of Space and Time (LSST) conducted by Vera C. Rubin Observatory in combination with space telescopes such as HST, James Webb Space Telescope, and Nancy Grace Roman Space Telescope. These results suggest that Neo could be used to improve both cosmological and galaxy evolution analyses based on massive, ground-based survey datasets like LSST. The model code is open source and available at https://purl.archive.org/neo/code.
\subsubsection*{Selected figures}
\begin{figure}[H]
\centering
\includegraphics[width=\linewidth,height=0.42\textheight,keepaspectratio]{/home/kf/research-assistant/data/cache/figures/http-arxiv.org-abs-2604.20195v1/figure\_1.png}
\captionsetup{labelformat=empty}\caption{Figure 1. Schematic of the Neo generator network including pre-processing and post-processing steps. In the Neo neural
network, the generator is defined by a contraction phase, expansion phase, and super-resolution phase. The contraction operators
ψf,b reduce the spatial extent of the tensors by a f}
\end{figure}
\newpage
\subsection*{9. Photon and neutrino fluxes from spheroidal dwarf galaxies in a decaying DM model}
\textbf{Focus:} galaxies, dark matter\\
\textbf{Authors:} A. Carrillo-Monteverde, L. López-Lozano, F. San Juan-Villegas\\
\textbf{Published:} 2026-04-22T01:06:35+00:00\\
\textbf{PDF:} \href{https://arxiv.org/pdf/2604.20086v1}{https://arxiv.org/pdf/2604.20086v1}\\
\textbf{Summary:}
In this work, we investigate a decaying dark matter scenario and its associated indirect detection signatures. The model consists of a scalar singlet with a lifetime exceeding the age of the Universe. Stability is ensured by a $Z_2$ symmetry imposed on the Lagrangian, allowing decay through a non-minimal gravitational coupling. The decay of dark matter produces Standard Model particles, which subsequently yield products such as gamma rays, neutrinos, and charged particles. We computed the gamma-ray and neutrino fluxes generated by this candidate in the Milky Way and in 14 dwarf spheroidal galaxies, as well as the corresponding expected number of events in selected experiments, using dedicated numerical tools. Results are presented for three benchmark masses and three coupling values consistent with cosmological constraints, showing that the predicted signals can be observable in specific regions of parameter space.
\subsubsection*{Selected figures}
\begin{figure}[H]
\centering
\includegraphics[width=\linewidth,height=0.42\textheight,keepaspectratio]{/home/kf/research-assistant/data/cache/figures/http-arxiv.org-abs-2604.20086v1/figure\_1.png}
\captionsetup{labelformat=empty}\caption{ve of D-factors from the Milky Way and som
onds to the D-factor of the Milky Way and th
The order of the magnitude is similar, so the
much from the dSphs galaxies. he objects studied in this work. The results}
\end{figure}
\newpage
\subsection*{10. Cold molecular gas distribution and kinematics in the low-metallicity dusty starburst of Mrk 996 resolved with ALMA}
\textbf{Focus:} galaxies\\
\textbf{Authors:} R. Slater, R. Amorín, J. A. Fernández-Ontiveros, F. J. Sáez-Ruiz, M. S. Oey, B. L. James, M. Mingozzi, M. Llerena, M. G. del Valle-Espinosa, K. Harrington, N. Kumari, R. Sánchez-Janssen, J. M. Vílchez\\
\textbf{Published:} 2026-04-22T00:54:21+00:00\\
\textbf{PDF:} \href{https://arxiv.org/pdf/2604.20080v1}{https://arxiv.org/pdf/2604.20080v1}\\
\textbf{Summary:}
Detecting cold molecular gas in metal-poor starbursts remains a major challenge. Low carbon and oxygen abundances hinder CO formation, while low dust content reduces shielding against UV photodissociation. Consequently, CO, the main tracer of molecular gas, becomes faint or undetectable. We study the spatial distribution and kinematics of cold molecular gas in Mrk 996, a nearby low-mass Wolf-Rayet galaxy hosting a dense, low-metallicity (about 1/5 solar) and nitrogen-enriched nuclear starburst with complex ionized gas kinematics. Using ALMA observations of CO(1-0) and CO(2-1), we map the morphology and kinematics of the molecular gas and compare them with optical and UV data, tracing the ionized gas and young stellar populations. We detect compact CO clouds within 800 pc of the starburst, spatially offset from the nuclear super star cluster (SSC) and the most highly ionized regions. The CO lines are narrow and supersonic, exhibiting velocity gradients with a mild global blueshift, indicating dynamically perturbed gas without evidence for fast outflows, in contrast with the highly ionized phase. The global CO(2-1)/CO(1-0) ratio is low (R21 \textasciitilde{} 0.3), consistent with subthermal excitation. The millimeter continuum peaks at the SSC, while CO emission is displaced toward obscured regions, suggesting it traces dense shielded clumps. ALMA recovers about half of the single-dish flux, indicating the presence of extended, low-surface-brightness molecular gas. Using a metallicity-dependent CO-to-H2 conversion factor, we infer a molecular gas mass of a few 10\textasciicircum{}7 solar masses. The molecular gas is only weakly coupled to the stellar feedback that dominates the ionized phase. Our results support a multiphase scenario in which dense molecular clumps survive in shielded regions, while CO is photodissociated in their envelopes, leaving a significant CO-dark H2 component (Abridged).
\subsubsection*{Selected figures}
\begin{figure}[H]
\centering
\includegraphics[width=\linewidth,height=0.42\textheight,keepaspectratio]{/home/kf/research-assistant/data/cache/figures/http-arxiv.org-abs-2604.20080v1/figure\_1.png}
\captionsetup{labelformat=empty}\caption{. 5. Optical view of Mrk 996: (Left) HST V+I map and (Right) V/I color map. CO contours are overlaid in the top panels, while dust-domina
tinuum contours at 115 GHz (band-3, red) and 230 GHz (band-6, green) are shown in the bottom panels. Contour levels are similar to th
Fig. 4. Coloured scale bars }
\end{figure}
\newpage
\subsection*{11. Purely Quadratic Non-Gaussianity from Tachyonic Instability: Primordial Black Holes and Scalar-Induced Gravitational Waves}
\textbf{Focus:} dark matter\\
\textbf{Authors:} He-Xu Zhang, Mei Huang\\
\textbf{Published:} 2026-04-21T23:58:51+00:00\\
\textbf{PDF:} \href{https://arxiv.org/pdf/2604.20063v1}{https://arxiv.org/pdf/2604.20063v1}\\
\textbf{Summary:}
We investigate primordial black hole (PBH) formation in a cosmological scenario where curvature perturbations follow purely quadratic non-Gaussianity, $ζ= A(\phi ^2-\langle\phi ^2\rangle)$, arising from tachyonic instability in multi-component inflationary models. Within an extended Press-Schechter framework based on the compaction function, we derive the probability distribution of the linear compaction function and its asymptotic exponential tail, demonstrating that PBH abundance is exponentially sensitive not only to the amplitude of perturbations but also to the correlation coefficient $ρ$ between the smoothed field and its radial gradient. We further find that, for $A<0$, the spectral width of the curvature power spectrum plays a decisive role in avoiding PBH overproduction: broad spectra yield mildly negative $ρ$ and fail to suppress PBH formation, while sufficiently narrow spectra drive $ρ\to -1$, resulting in an exponential suppression while maintaining a sizable gravitational-wave signal. Thermal inflation provides a useful benchmark scenario with asteroid-mass PBH dark matter and high-frequency scalar-induced gravitational waves potentially detectable by future space-based interferometers, but its typically broad spectra make it challenging to reconcile PTA observations with PBH constraints.
\newpage
\subsection*{12. Unraveling Chemical Enrichment in Extreme Emission-Line Galaxies: A Multi-Element Bayesian View of Bursty Star Formation and Galaxy Evolution in DESI}
\textbf{Focus:} galaxies\\
\textbf{Authors:} Razieh Emami, James A. A. Trussler, Tiger Yu-Yang Hsiao, Kaley Brauer, Lars Hernquist, Randall Smith, Douglas Finkbeiner, Fengwu Sun, Rebecca Davies, James F. Steiner, Mark Vogelsberger, Tobias Looser, Grant Tremblay, Letizia Bugiani\\
\textbf{Published:} 2026-04-21T23:42:10+00:00\\
\textbf{PDF:} \href{https://arxiv.org/pdf/2604.20060v1}{https://arxiv.org/pdf/2604.20060v1}\\
\textbf{Summary:}
Extreme emission-line galaxies (EELGs) probe chemical enrichment in low-mass, bursty systems where star formation, feedback, and gas accretion are poorly constrained. Using DESI DR1, we select 23 nearby EELGs with detections of 19 ionic species (S/N $\geq$ 4), stellar masses $ M_* \geq 10^7 M_{\odot}$, and extreme H$\alpha $ and [O III] 5007 equivalent widths (EW $\geq$ 500 Angstrom). We infer non-parametric star-formation histories and fit a Bayesian single-zone chemical-evolution model to O, N, Ne, S, and Ar, allowing time-dependent star-formation efficiency, outflow mass loading, and evolving inflow metallicity. We find short depletion timescales and large mass-loading factors, indicating rapid gas cycling in a burst-driven, non-equilibrium regime, with depletion times below Kennicutt-Schmidt expectations. Star-formation efficiency and outflows are well constrained, while inflow metallicity is weaker due to degeneracies with metal production. Abundance ratios isolate physical drivers: star-formation efficiency sets evolutionary tracks, outflows regulate metal retention and X/O normalization, and inflow metallicity sets baseline enrichment. N/O strongly constrains burst timing and gas flows, Ne/O remains nearly invariant, and S/O and Ar/O show intermediate sensitivity. These results demonstrate that multi-element abundances provide a direct probe of baryon-cycle processes in extreme low-mass starbursts.
\subsubsection*{Selected figures}
\begin{figure}[H]
\centering
\includegraphics[width=\linewidth,height=0.42\textheight,keepaspectratio]{/home/kf/research-assistant/data/cache/figures/http-arxiv.org-abs-2604.20060v1/figure\_1.png}
\captionsetup{labelformat=empty}\caption{Figure 8. Same as Figure 7, but for the star formation efficiency scaling parameters 𝑛𝜏and log10(𝐴𝜏). The inferred values of 𝑛𝜏exhibit a high
degree of uniformity across the sample, with strong overlap between different SFH realizations and metallicity assumptions. In contrast, the
normalization log}
\end{figure}
\newpage
\subsection*{13. MAUVE-MUSE: Ionization and Kinematic Signatures of Environmental Effects on Virgo Cluster Disks}
\textbf{Focus:} galaxies\\
\textbf{Authors:} Toby Brown, Luca Cortese, Barbara Catinella, A. Fraser-McKelvie, Adam B. Watts, Amirnezam Amiri, Alessandro Boselli, Woorak Choi, Aeree Chung, Timothy A. Davis, Eric Emsellem, Pavel Jáchym, María J. Jiménez-Donaire, Tutku Kolcu, Bumhyun Lee, Andrei Ristea, Jesse van de Sande, Kristine Spekkens, Sabine Thater, Christine D. Wilson, Nikki Zabel\\
\textbf{Published:} 2026-04-21T23:32:01+00:00\\
\textbf{PDF:} \href{https://arxiv.org/pdf/2604.20056v1}{https://arxiv.org/pdf/2604.20056v1}\\
\textbf{Summary:}
We present early science results from the MAUVE (Multiphase Astrophysics to Unveil the Virgo Environment) program which targets 40 Virgo Cluster galaxies to investigate the effect of environment on the interstellar medium (ISM) at \textasciitilde{}100 pc scales. From 12 galaxies in the MAUVE-MUSE early sample, we find systematically elevated line ratios compared to PHANGS-MUSE field disks, with higher medians of [N II]/H$\alpha $ (0.75 vs. 0.50), [S II]/H$\alpha $ (0.57 vs. 0.49), and [O III]/H$\beta $ (1.04 vs. 0.68). Spatially resolved BPT diagrams show 74\% of MAUVE-MUSE spaxels ionized by sources other than H II regions, versus 61\% in the field, and we find these ionization differences to be closely coupled to broadened kinematics. 44\% of MAUVE-MUSE spaxels exceed H$\alpha $ $\sigma _{LOS} = 40$ km/s (vs. 26\% in the field), driven mainly by non-star-forming gas with $\sigma _{LOS}$ between 40 and 80 km/s, consistent with enhanced contribution of diffuse ionized gas (DIG). A subdominant tail of 5\% of spaxels at $\sigma _{LOS} > 100$ km/s, largely absent in PHANGS-MUSE (1\%), points to shocks or turbulent mixing layers from intracluster interactions. Our results show that environmental quenching primarily suppresses star formation, unveiling DIG as the dominant ionized component in cluster disks. The elevated line ratios and broadened kinematics observed in the MAUVE sample reflect the physical state of the ISM in the absence of vigorous star formation, rather than widespread direct environmental excitation. The observed shock-like emission provides an additional, secondary contribution likely driven by active interactions with the intracluster medium.
\subsubsection*{Selected figures}
\begin{figure}[H]
\centering
\includegraphics[width=\linewidth,height=0.42\textheight,keepaspectratio]{/home/kf/research-assistant/data/cache/figures/http-arxiv.org-abs-2604.20056v1/figure\_1.png}
\captionsetup{labelformat=empty}\caption{tially resolved BPT diagrams for the PHANGS--MUSE (top) and MAUVE--MUSE (bottom
i]/H$\textbackslash{}beta$. Right: [S ii]/H$\textbackslash{}alpha$ vs. [O iii]/H$\textbackslash{}beta$. Colors indicate the orthogonal distance, dSF, in dex fr
nd theoretical Ke01 (solid) star-formation boundaries, where positive values denote increasin
on excitation toward non--H ii re}
\end{figure}
\end{document}
